\chapter{State of the Art}
\label{cap:stato_arte}

This chapter reviews vision-based methods for agricultural row
navigation, focusing on the conversion of image observations into a
local guidance reference. It covers visual perception, geometric
extraction, and the integration of these stages with robot control.
Dataset preparation and the implementation used in this thesis are
described in Chapters~\ref{cap:dataset_modelli}
and~\ref{cap:geometric_reference}.

\section{Agricultural row navigation and sensing}
\label{sec:navigazione_contesto}

\subsection{Navigation tasks and operating conditions}
\label{subsec:scenari_navigazione}

Vision-based agricultural navigation uses image features to identify
crop rows, crop boundaries, or the soil region between rows, and then
converts them into a reference for the robot \cite{bai2023vision}.
These are related but distinct tasks: following the center of a crop
row, staying in the inter-row path, and detecting the boundary between
standing crop and harvested ground do not produce the same visual
reference. The reviewed studies address greenhouse cucumber rows,
sugar-beet fields, vineyards, orchards, and wheat or rapeseed fields
\cite{chen2021navigation,silva2023rowdetection,mendez2023vineyards,
navone2025gpsfree,kneip2020cropedge}.

Field appearance varies with illumination, shadows, weeds, plant size,
occlusion, and gaps in the rows. These conditions affect both
segmentation and the geometric reference derived from it
\cite{silva2023rowdetection,navone2025gpsfree}.

\subsection{Sensing for local guidance}
\label{subsec:vision_based}

GNSS can provide global position outdoors, but crop canopies may
obstruct satellite signals. LiDAR, stereo, and RGB-D cameras provide
depth or three-dimensional structure; their suitability depends on the
field environment, sensor configuration, and cost. A monocular camera
is less complex and provides appearance information, but does not
directly measure depth \cite[Sections~2.1--2.3]{bai2023vision}.
Navone et al. use RGB-D sensing for local guidance in dense rows where
vegetation can also obstruct GNSS reception \cite{navone2025gpsfree}.

This thesis focuses on camera-based local guidance. It does not address
global localization or reconstruct the field in three dimensions.

\section{Visual perception of crops and inter-row paths}
\label{sec:percezione_sentiero}

\subsection{Classical image-processing methods}
\label{subsec:metodi_classici_colore}

Classical pipelines combine image filtering, thresholding, and
application-specific geometric rules. Mendez et al. convert vineyard
images to grayscale, apply Gaussian filtering and Otsu thresholding,
select the largest blob, and use column pixel counts to estimate the row
center \cite[Sections~3.4--4]{mendez2023vineyards}. Such methods can be
effective in the conditions for which their parameters were selected,
but changes in illumination, weeds, and crop appearance can require
retuning \cite[Section~2]{kneip2020cropedge}.

\subsection{Learning-based segmentation}
\label{subsec:deep_learning_segmentazione}

Semantic segmentation assigns a class to each pixel, but the label
definition determines what the geometric stage can recover. De Silva
et al. train U-Net to predict narrow crop-row masks in sugar-beet
images; the labels represent row structure rather than individual plant
pixels or traversable soil \cite[Sections~3.3 and~4.1]{silva2023rowdetection}.
Navone et al. instead predict a binary vegetation mask from RGB images
using a MobileNetV3 backbone and an LR-ASPP head, then combine that
output with depth for reference extraction
\cite[Sections~3.3 and~4.1]{navone2025gpsfree}.

The distinction matters for this work: its target is the soil region
between crop rows, not a crop-row line or a vegetation mask. A
segmentation model therefore supplies a candidate region; the geometric
stage must still select the corridor and construct the guidance
reference.

\subsection{Training data and model-assisted annotation}
\label{sec:dataset_transfer}

Pixel-level annotation is time-consuming, and agricultural datasets
must represent changes in crop, growth stage, and lighting. De Silva et
al. include 11 field variations in their dataset
\cite[Section~3]{silva2023rowdetection}. Navone et al. initialize their
segmentation model with ImageNet weights and train using synthetic
AgriSeg images supplemented by 100 real crop images
\cite[Section~4.1, Table~1]{navone2025gpsfree}. Pretraining and the
choice of training images are separate design choices.

Segment Anything Model~3 (SAM~3) supports text- and image-prompted
concept segmentation. Its training data were assembled using a data
engine with human and automated quality checks
\cite[Sections~2 and~4]{carion2026sam3}. In this thesis, SAM~3 is used
to prepare labels offline, not as the network running during robot
navigation. The quality of these masks and the time required to review
or correct them are evaluated separately in Chapter~\ref{cap:dataset_modelli}.

\section{From segmentation to a navigation reference}
\label{sec:estrazione_traiettoria}

The geometric stage reduces a mask to a quantity the controller can use.
Depending on the representation and method, the output may be a
horizontal image position, a straight line, a crop-edge estimate, or a
curve. These outputs describe different tasks and coordinate systems.

\subsection{Line fitting and crop-edge detection}
\label{subsec:hough_fitting}

Chen et al. estimate a path line for greenhouse cucumber rows from
navigation points selected in plant--soil segmentation. Their
prediction-point Hough transform uses a least-squares estimate to
restrict the line search and is compared with traditional Hough and
least-squares fitting \cite[Sections~2--3]{chen2021navigation}.

Kneip et al. detect crop--ground transitions in scanlines of a stereo
elevation map, then fit a linear crop-edge model using a Huber
M-estimator and iteratively reweighted least squares. Their output is a
crop boundary in ground-plane coordinates, not an image-space estimate
of the inter-row path \cite[Sections~6--7]{kneip2020cropedge}.

\subsection{Column-wise histogram methods}
\label{subsec:metodi_locali}

Histogram methods aggregate mask pixels by image column. Mendez et al.
estimate the row center from the five columns with the largest pixel
counts in a thresholded background region
\cite[Section~3.4.3]{mendez2023vineyards}. Navone et al. use a different
input: a temporally combined vegetation mask refined with depth. SegMin
smooths its column histogram and selects the minimum; SegMinD weights
the vegetation mask by normalized inverse depth before calculating the
histogram \cite[Section~3]{navone2025gpsfree}. In both cases the output
is a horizontal position, not a full path line.

\subsection{Horizontal strips and adaptive multi-ROI methods}
\label{subsec:strips}

Li et al. combine a Faster-U-net inter-row-path mask with an adaptive
multi-ROI procedure. Horizontal regions supply ordered local points,
which are used as control points for a multisegment cubic B-spline.
Their experiments use 40-pixel search regions and discuss the trade-off
between the number of points and fitting behavior
\cite[Section~2.2.3, Figure~9]{li2022multiroi}.

This is the closest geometric comparison for the present work because
both pipelines select local evidence from horizontal image regions.
The representation and fitted model differ: Li et al. use an
inter-row-path mask and a spline, whereas this thesis estimates a local
straight reference from the soil corridor between parallel crop rows.

\subsection{Triangle Scan Method}
\label{subsec:triangle_scan}

De Silva et al. propose the Triangle Scan Method (TSM) for crop-row
masks. It first scans an upper horizontal band to find an anchor column,
using a preset anchor if the scan fails a threshold. Candidate lines
from that anchor to positions along the base of a triangular region are
then evaluated; the line with the greatest mask support defines the
central crop row. Its two endpoints provide position and angular errors
for visual servoing \cite[Section~4.2, Algorithm~1]{silva2023rowdetection}.

TSM is not the same as extracting one point from each of several
horizontal strips: the upper scan establishes an anchor for a bounded
line search. It also operates on crop-row masks, so its output is not
directly equivalent to a reference extracted from a soil mask.

\subsection{Comparison of geometric representations}
\label{subsec:confronto_riferimenti}

Table~\ref{tab:confronto_riferimenti_geometrici} summarizes the input
representation and geometric output of the methods discussed above.

\begin{table}[htbp]
    \centering
    \footnotesize
    \setlength{\tabcolsep}{3pt}
    \renewcommand{\arraystretch}{1.15}
    \caption{Perceptual inputs and geometric references in selected
    agricultural navigation studies.}
    \label{tab:confronto_riferimenti_geometrici}
    \begin{tabularx}{\textwidth}{@{}>{\raggedright\arraybackslash}p{0.16\textwidth}
        >{\raggedright\arraybackslash}X
        >{\raggedright\arraybackslash}X
        >{\raggedright\arraybackslash}p{0.22\textwidth}@{}}
        \toprule
        Study & Perceptual input & Geometric operation & Output \\
        \midrule
        Chen et al.\ \cite{chen2021navigation}
        & Plant--soil segmentation and navigation points
        & Restricted Hough search around a predicted point
        & Image-space path line \\
        \addlinespace
        Kneip et al.\ \cite{kneip2020cropedge}
        & Stereo elevation map
        & Scanline transitions and robust linear fit
        & Ground-plane crop edge \\
        \addlinespace
        Mendez et al.\ \cite{mendez2023vineyards}
        & Thresholded background region
        & Column counts and mean of the five largest
        & Horizontal row-center position \\
        \addlinespace
        Navone et al.\ \cite{navone2025gpsfree}
        & Vegetation mask with depth
        & Smoothed column histogram; inverse-depth weighting in SegMinD
        & Histogram-minimum position \\
        \addlinespace
        Li et al.\ \cite{li2022multiroi}
        & Predicted inter-row-path mask
        & Adaptive horizontal regions and cubic B-spline
        & Curved image-space path \\
        \addlinespace
        De Silva et al.\ \cite{silva2023rowdetection}
        & Predicted crop-row mask
        & Upper anchor and triangular line search
        & Central crop-row line \\
        \bottomrule
    \end{tabularx}
\end{table}

The methods are not directly ranked by this comparison: their inputs
and outputs differ. Li et al. provide the closest structural precedent
for strip-based point selection, while TSM illustrates a line search
for a crop-row representation. Neither method uses the same
soil-corridor mask and straight-line model as the pipeline developed
in this thesis.

\section{System integration and experimental evidence}
\label{sec:sistemi_completi}

\subsection{Perception and control}
\label{subsec:architetture_percezione_controllo}

Navigation systems use the extracted reference to generate motion
commands. Navone et al. convert the displacement of their histogram
reference into linear and angular velocity commands and smooth them
with an Exponential Moving Average
\cite[Section~3.4]{navone2025gpsfree}. Mendez et al. use proportional
control based on row-center displacement and the distribution of
background pixels \cite[Section~3.5]{mendez2023vineyards}. De Silva et
al. combine the central row's angular and positional errors in a
proportional visual-servoing controller for simulation
\cite[Section~4.3]{silva2023rowdetection}.

\subsection{Scope of reported evaluations}
\label{subsec:validation_scope}

The studies evaluate different stages of navigation. Kneip et al.
assess crop-edge detection on recorded harvesting data
\cite[Section~8]{kneip2020cropedge}; Chen et al. report angular errors
and processing times for image-based path fitting
\cite[Section~3.6, Tables~6--7]{chen2021navigation}. De Silva et al.
evaluate row detection on real images and visual servoing in simulation,
with preliminary real-world row following also reported
\cite[Sections~5.2--5.3 and~6, Table~5]{silva2023rowdetection}.
Mendez et al. report closed-loop vineyard trials, measuring centering
error from the distances to the two sides of the row
\cite[Sections~3.6 and~4]{mendez2023vineyards}. Navone et al. test in
simulation and crop fields, using LiDAR and attitude measurements to
assess displacement and heading \cite[Sections~4.3--4.4]{navone2025gpsfree}.

Segmentation IoU, image-space line error, field tracking error, and
processing time measure different parts of the system. Results from
these studies therefore provide context, not a common performance
benchmark.

\section{Positioning of this thesis}
\label{sec:considerazioni_finali_stato_arte}

This thesis studies a camera-based pipeline that segments plants and
the soil corridor, selects the relevant inter-row region, and estimates
a local guidance line for navigation between parallel crop rows. Its
geometric stage uses horizontal strips to select candidate intervals
and fits the two corridor boundaries separately. This shares the
local-selection structure of Li et al.'s multi-ROI method, but uses a
straight image-space reference rather than a spline. The experiments
in the following chapters evaluate the segmentation models, geometric
extractors, and integrated robot at their respective stages.